\begin{claim}{C.30}{Proved}{Properness of a transitive-action centralizer}
Let $V$ be a connected second-countable surface without boundary,
and let $K$ be a closed subgroup of $\operatorname{Homeo}(V)$ acting
transitively on $V$, with the compact-open topology. Put
$N=C_{\operatorname{Homeo}(V)}(K)$. Then $N$ is locally compact,
its action on $V$ is free and proper, and for every $x\in V$
evaluation is a homeomorphism onto the closed subset
\[
 N\xrightarrow{\ \mathrm{ev}_x\ }
 N_K(K_x)/K_x\ \subseteq\ K/K_x\cong V.                  \tag{C.30.1}
\]
Under the normalizer-quotient group structure, this is an
anti-isomorphism of topological groups. For an open orbit $O$ of the
loop-generated group in C.28, the image in $\operatorname{Homeo}(O)$
of its operator stabilizer lies in such an $N$; the closure of that
image acts freely and properly on $O$.
\end{claim}
\begin{proof}
The group $K$ is Polish. The open mapping theorem
\cite[Theorem~3.5]{C_GomezPinedoUzcategui2017}, with the ordinary
continuous transitive action on the nonmeagre Hausdorff surface $V$,
identifies $K/K_x$ homeomorphically with $V$. In particular the
orbit map $k\mapsto kx$ is open. We will use the following exact
consequence, requiring no local continuous section: for every
sequence $x_j\to x$, there are $k_j\to1$ in $K$ with $k_jx=x_j$.
To prove it, choose a decreasing identity-neighborhood basis $W_m$.
The open sets $W_mx$ contain $x$, so successive tails of the sequence
lie in them. Choose $k_j$ from the corresponding $W_{m(j)}$, with
$m(j)\to\infty$.

An element of $N$ is determined by its value at $x$, since
$n(kx)=kn(x)$. If $n(x)=ax$, equality of stabilizers gives
$aK_xa^{-1}=K_x$. Conversely, for $a$ in this normalizer the map
$kK_x\mapsto kaK_x$ is a homeomorphism commuting with $K$.
The normalizer is closed: both inclusions defining
$aK_xa^{-1}=K_x$ are intersections of closed conditions.
It is saturated for the quotient by $K_x$, so its image is a closed
subset of $K/K_x=V$. This proves the claimed set-theoretic image
and closedness in (C.30.1). If $n(x)=ax$ and $m(x)=bx$, then
$(nm)(x)=bax$, proving the stated anti-isomorphism.

We prove inverse continuity of evaluation, including the issue of
noncompact $V$. Suppose that $n_j(x)\to n(x)$, with $n_j,n\in N$.
For any compact set $B\subset V$, a failure of uniform convergence
on $B$ would provide a subsequence and $y_j\in B$ with
$y_j\to y$ witnessing that failure. Choose $a\in K$ with $ax=y$,
and then $b_j\to1$ with $b_jy=y_j$, using open evaluation at $y$.
Commutation gives
\[
 n_j(y_j)=n_j(b_jax)=b_ja\,n_j(x)\longrightarrow a n(x)=n(y).
                                                               \tag{C.30.2}
\]
Also $n(y_j)\to n(y)$, a contradiction. Hence $n_j\to n$ uniformly
on every compact set, which is compact-open convergence.
Inversion can be checked directly as well. Put $c_j=n^{-1}n_j$,
so $c_jx\to x$, and choose $a_j\to1$ with $a_jx=c_jx$.
Because $c_j$ commutes with $a_j$,
$c_j^{-1}x=a_j^{-1}x\to x$; applying (C.30.2) to $c_j^{-1}$
gives convergence of the inverses. All the spaces are metrizable,
so this sequential argument proves inverse continuity.
Thus (C.30.1) is a homeomorphism, and $N$ is locally compact as a
closed subspace of the locally compact surface.

If $n\in N$ fixes one point, transitivity and commutation make it
fix every point, proving freeness. To prove properness, consider
sequences with
\[
                  x_j\longrightarrow x,\qquad
                  n_jx_j\longrightarrow y.
\]
Choose $k_j\to1$ with $k_jx=x_j$. Then
\[
                n_jx=k_j^{-1}n_jx_j\longrightarrow y.    \tag{C.30.3}
\]
The image of evaluation at $x$ is closed, so $y=nx$ for a unique
$n\in N$. Its inverse continuity gives $n_j\to n$.
For any compact subset of $V\times V$, every sequence in its
inverse image under $(n,z)\mapsto(z,nz)$ therefore has a convergent
subsequence, with limit in the same inverse image. Metrizability
turns this sequential compactness into compactness, proving that
the action map is proper. Every closed subgroup of $N$ inherits
the free proper action.

Finally let $G$ be the group in C.28 and let $O$ be one of its open
orbits. It is connected because every element of $G$ is joined to
the identity by a path in $G$. The closure
\[
        K=\overline{\{g|_O:g\in G\}}
                       ^{\operatorname{Homeo}(O)}
\]
is a closed Polish group still transitive on $O$. Let $\widetilde H$
be any group commuting with $G$, and let $\widetilde H_O$ be its
setwise stabilizer of $O$. The precise restriction homomorphism is
\[
 \widetilde H_O\longrightarrow\operatorname{Homeo}(O),\qquad
 h\longmapsto h|_O,
 \quad
 \ker=\{h\in\widetilde H_O:h|_O=\id\}.                 \tag{C.30.4}
\]
Its image commutes with the restriction of $G$ and hence with $K$.
The closure of this image is a closed subgroup of $N$, and therefore
acts freely and properly on $O$, as claimed.
\end{proof}

\paragraph{Precise remaining use.}
The closure in the last assertion has its subspace topology from
$\operatorname{Homeo}(O)$. Formula (C.30.4) does not assert that the
original discrete operator group acts freely, since its kernel may
be nontrivial; nor does it equip a possibly dense abstract image with
the discrete topology and claim properness for that topology.
The result supplies the closed locally compact transformation group
actually present on an open orbit. No Lie-group classification or
quotient-bundle contraction is invoked here. A contraction relative
to orbit frontiers, and its continuous assembly with C.27 and C.29
on the one-dimensional strata, remain unproved. Thus C.30 does not
change the unresolved status of C.12.
